


The Spanish Civil War starts in 1936 and opposes the Republican forces - composed of left-wing factions as the communists, anarchists, and socialists - and the Nationalists groups, led by General Francisco Franco and supported by fascist and conservative groups. The Nationalists won in 1939, establishing a dictatorship that lasted until 1977. By early 1939, Catalonia was one of the last Republican territories and had welcomed many internally displaced people. After the fall of Barcelona and the Nationalist offensive across Catalonia, the Republican defeat became inevitable. This collapse triggered the mass exodus of about 500,000 Spanish Republicans to France between late January and mid-February 1939. Both combatants and civilians crossed the Pyrenees in search of refuge, producing the largest sudden population movement ever recorded at the Franco-Spanish border. This influx represented a substantial demographic shock for France, representing 22\% of the country’s total foreign population according to the last pre-war census of 1939.

%Considering the last census available of 1936, the total population of France was 41967056 habitants, this represents 1,19% of the population. The shock represents 22% of the total foreign population in 1939. The total foreign population was 2198236 in 1939. The most important groups are: 721000 italians which represented 33% of the foreign population. Polish: 423000 (19percent) and Spanish 254000 (12percent). 

This sudden influx, concentrated over a few weeks, overwhelmed the French authorities who were unprepared for its scale \citep{rafaneau20081939}. Worn down by economic crisis, plagued by xenophobic sentiment, and increasingly isolationist, French authorities offers Spanish refugees a hostile reception \citep{Coignard2012}. For French authorities, still reeling from World War I, the priority was to maintain the public order and security. They were particularly worried about the spread of radical left-wing ideas, so the incoming Spaniards were treated as a security problem and were interned and dispersed all over the country \citep{rafaneau20081939}. 

Families were often separated at the border, and the civilian population was transferred to reception centres dispersed across the country. Around 170,000 of them were women, children, and elderly people \citep{perez_rodriguez_indeseables_2022}. Meanwhile, most men (particularly those of military age) were interned in improvised camps - ``concentration camps” as they were officially called - near the border in the south \citep{DreyfusArmand1999}. The legal basis for this camp system had been laid shortly before: in November 1938, the Daladier government passed a decree-law allowing the internment of “undesirable foreigners” under surveillance. The Spanish were the first to suffer the consequences of this new policy \citep{eurom_argeles-sur-mer_2014}. 

As no permanent facilities had been prepared in advance, the selection of camp sites was improvised. The very first camps were in the Pyrénées-Orientales (French Catalonia) along the Mediterranean coast: notably Argelès-sur-Mer, Saint-Cyprien, and Le Barcarès, all essentially stretches of beach beach and open fields demarcated and enclosed with barbed wire \citep{Coignard2012}. Soon, additional camps were created deeper inside France to alleviate the severe overcrowding. By the spring of 1939, a second ring of internment camps extended across the Pyrenees and southwestern France. Key camps included Gurs (in Béarn, near Pau), Le Vernet d’Ariège (in the Ariège department), Septfonds (Tarn-et-Garonne), Bram (Aude), Agde (Hérault), and Rieucros (Lozère) \citep{Coignard2012}. In the rest of the departments, the types of accommodations varied widely, whether suitable for hosting people (schools, holiday centres, sanatoriums, barracks, uninhabited houses...) or not initially appropriate (former factories, stud farms, stables, mills, village halls) \citep{DreyfusArmand1999}.

%if it was the prefect who had to give the census of places, they are elected each two years and not on the same level than the commune. so you can argue that the political tendencies did not matter that much. 

From the moment the refugees arrived, French authorities aimed to reduce the number of Spaniards on French soil. Concerned about the financial cost, they  promote ``voluntary" return through incentives and coercion. By early August, around 250,000 refugees had returned to Spain \citep{perez_rodriguez_indeseables_2022}. Moreover, around 20,000 people emigrated to other countries, with over 15,000 going to Latin America. As a result, the number of Spanish refugees and political exiles in French territory had reached approximately 180,000 individuals by December 1939 \citep{rafaneau20081939}. As World War II approached, France began to see those who remained as a valuable labour force for its war effort. In April 12, 1939, Foreign Worker Companies (Compagnies de Travailleurs Étrangers, CTE) were established to conscript refugees into labor units \citep{Coignard2012}. By the end of December, 50,000 to 60,000 Spanish refugees were working in these groups particularly along the Italian border and on Maginot Line fortifications \citep{rafaneau20081939}. In addition, the French military encouraged Spanish veterans to enlist in the armed forces. Thousands chose to join the French Foreign Legion or other military groups as a way to leave the camps \citep{perez_rodriguez_indeseables_2022}. When France fell to German forces in June 1940, Spanish refugees faced new dangers. Many men were captured and held as prisoners of war, and thousands of Spanish Republicans were subsequently deported to Nazi concentration camps—most notoriously Mauthausen, where around 7,000 died. Others managed to evade capture and join the French Resistance in significant numbers. It is estimated that up to 15,000 Spaniards participated in the French Resistance or served in the Free French forces between 1940 and 1944 \citep{academicoup}. By the end of the Second World War, it is estimated that 100000 Spanish republicans were still in the country \citep{perez_rodriguez_indeseables_2022}













